\chapter{微分の応用}\label{章:微分の応用}
微分すると,一点での変化率が得られた.その値を区間全体で調べれば,増減や最大値も分かるだろうか.まず,一点での傾きと二点間の傾きを比べる平均値の定理を考えよう.
\section{平均値の定理}\label{節:平均値定理}
微分係数は1点での変化率である.一方,関数が区間全体で増えるかどうかを知るには,離れた2点の値を比べる必要がある.では,1点での傾きから,2点での値の差について何が言えるだろうか.

まずは簡単な関数で実験してみよう.区間の両端を結ぶ直線(割線)の傾きを求め,接線の傾き $f'(c)$ がそれに等しくなる点 $c$ を区間の内側で探すと,次のようになる.
\begin{center}
\begin{tabular}{cccc}
\toprule
$f(x)$ & 区間 & 割線の傾き & $c$\\
\midrule
$x^2$ & $[0,2]$ & $2$ & $1$\\
$x^3$ & $[0,3]$ & $9$ & $\sqrt3\approx1.73$\\
$x^3-x$ & $[-2,2]$ & $3$ & $\pm\frac{2}{\sqrt3}\approx\pm1.15$\\
\bottomrule
\end{tabular}
\end{center}
たとえば2行目では,割線の傾きは $\frac{27-0}{3-0}=9$ であり,\,$3c^2=9$ を解いて $c=\sqrt3$ を得る.図\ref{図:平均値定理}は1行目の場合を描いたものである.

\begin{figure}[htbp]
\centering
\begin{tikzpicture}[x=2.6cm,y=1.1cm,>=Stealth]
 \draw[->] (-0.15,0)--(2.35,0) node[right]{$x$};
 \draw[->] (0,-0.4)--(0,4.7) node[above]{$y$};
 \draw[very thick,domain=0:2.12,samples=70] plot (\x,{\x*\x});
 \draw[blue,thick] (0,0)--(2,4);
 \draw[red,thick,domain=0.4:1.8] plot (\x,{2*\x-1});
 \fill (0,0) circle(1.6pt);
 \fill (2,4) circle(1.6pt) node[above left]{$(2,4)$};
 \fill (1,1) circle(1.6pt) node[below right]{$(1,1)$};
 \draw[dashed] (1,0) node[below]{$1$}--(1,1);
 \draw[dashed] (2,0) node[below]{$2$}--(2,4);
 \node[blue] at (0.63,2.05) {割線};
 \node[red] at (1.8,2.2) {接線};
 \node at (1.35,4.55) {$y=x^2$};
\end{tikzpicture}
\caption{平均変化率に等しい微分係数をもつ点が,区間の内側に存在する.}
\label{図:平均値定理}
\end{figure}

どの例でも,割線と平行な接線が区間の内側で引けている.点 $c$ は区間の中央とは限らず,3行目のように1つに決まるとも限らない.そこで,「割線と傾きが等しい接線が,区間の内側のどこかで引ける」と予想できる.これを証明するために,まず極値で何が起こるかを確かめよう.

点 $c$ の十分近くのすべての $x$ で $f(x)\le f(c)$ なら,\,$f$ は $c$ で\textbf{極大値}をとるという.不等号を逆にしたものが\textbf{極小値}である.極大・極小は近くの点との比較であり,区間全体での最大・最小とは区別する.

$c$ が区間の内点で,\,$f$ がそこで微分可能かつ極大であるとしよう.\,$|h|$ が小さいとき $f(c+h)-f(c)\le0$ だから,差商の符号は $h$ の符号で決まる.
\[
 \frac{f(c+h)-f(c)}{h}\le0\quad(h>0),\qquad
 \frac{f(c+h)-f(c)}{h}\ge0\quad(h<0).
\]
$f$ は $c$ で微分可能なので,\,$h\to0+$ と $h\to0-$ の極限はともに $f'(c)$ に等しい.したがって $f'(c)\le0$ かつ $f'(c)\ge0$ となり,\,$f'(c)=0$ なのだ.極小の場合も不等号が逆になるだけである.ただし,端点では片側からしか比較できないので,この結論は使えない.

予想のうち,割線が水平な場合,つまり $f(a)=f(b)$ の場合から考えよう.両端の高さが等しければ,グラフは途中で山か谷を作るだろう.その頂上や谷底では,先ほどの議論から接線が水平になるはずである.

\begin{theorem}{ロルの定理}\label{定理:ロル}
$a<b$ とする.\,$f$ が閉区間 $[a,b]$ で連続,開区間 $(a,b)$ で微分可能で,\,$f(a)=f(b)$ なら,
\[
 f'(c)=0
\]
を満たす $c\in(a,b)$ が少なくとも1つ存在する.
\end{theorem}

定理\ref{定理:連続関数存在}により,\,$f$ は $[a,b]$ で最大値と最小値をとる.\,$f$ が定数なら,どの内点でも $f'=0$ である.定数でなければ,最大値と最小値の少なくとも一方は端点の値 $f(a)=f(b)$ と異なる.その値をとる点 $c$ は内点であり,\,$f$ はそこで極値をとるから,\,$f'(c)=0$ となる.

割線が水平でない場合は,割線の傾きの分だけグラフを傾け直してやれば,水平な場合に帰着する.

\begin{theorem}{平均値の定理}\label{定理:平均値}
$a<b$ とする.\,$f$ が閉区間 $[a,b]$ で連続,開区間 $(a,b)$ で微分可能なら,
\[
 f'(c)=\frac{f(b)-f(a)}{b-a}
\]
を満たす $c\in(a,b)$ が少なくとも1つ存在する.
\end{theorem}

割線の傾きを $m=\frac{f(b)-f(a)}{b-a}$ と書き,
\[
 g(x)=f(x)-m(x-a)
\]
と置く.\,$g$ は $f$ から割線と同じ傾きの一次関数を引いたものであり,
\[
 g(a)=f(a),\qquad g(b)=f(b)-\{f(b)-f(a)\}=f(a)
\]
となって両端の値が等しい.\,$g$ は $[a,b]$ で連続,\,$(a,b)$ で微分可能だから,ロルの定理により,ある $c\in(a,b)$ で
\[
 g'(c)=f'(c)-m=0
\]
となる.よって $f'(c)=m$ であり,予想は正しかった.

条件の役割も確認しておこう.\,$f(x)=|x|$ を $[-1,1]$ で考えると,両端の値は等しいが,微分可能な点での導関数は $-1$ か $1$ であり,\,$0$ にはならない.\,$x=0$ で微分できないため,ロルの定理も平均値の定理も適用できないのである.導関数の連続性は定理の仮定に含まれていないが,区間の内側のすべての点で微分可能であるという仮定は外せない.

\subsection{導関数と増減}
平均値の定理を使えば,導関数の符号から関数の増減が分かる.まず言葉を決めておこう.区間内の任意の2点 $x_1<x_2$ についてつねに $f(x_1)<f(x_2)$ となるとき,\,$f$ はその区間で\textbf{狭義に増加する}(狭義単調増加)という.つねに $f(x_1)>f(x_2)$ となるときは\textbf{狭義に減少する}という.そのどちらかであることを\textbf{狭義単調}という.\label{用語:狭義単調}「狭義」は,等号 $f(x_1)=f(x_2)$ を許さないという意味である.

区間内の任意の2点 $x_1<x_2$ に平均値の定理を使うと,その間のある $c$ に対して
\[
 f(x_2)-f(x_1)=f'(c)(x_2-x_1)
\]
となる.\,$x_2-x_1>0$ だから,区間内で常に $f'>0$ なら $f(x_2)>f(x_1)$ であり,\,$f$ は狭義に増加する.\,$f'<0$ なら狭義に減少する.\,$f'\ge0$ の場合は $f(x_2)\ge f(x_1)$ までしか言えず,値が一定の部分もあり得る.

同じ議論から,次の結果も得られる.
\begin{theorem}{導関数が同じ関数の差}\label{定理:導関数同じ}
1つの区間のすべての内点で $F'=G'$ が成り立ち,端点を含む場合にはそこで両関数が連続なら,その区間で $F-G$ は定数である.特に,導関数が恒等的に $0$ なら,元の関数は定数である.
\end{theorem}

実際,\,$H=F-G$ と置くと $H'=0$ であり,上の式から任意の2点について $H(x_2)-H(x_1)=0$ となる.後の積分では,微分して同じ関数になる候補が,定数の違いを除いて同じであることをこの定理で確かめる.

ここで「1つの区間」という条件は外せない.たとえば
\[
 H(x)=
 \begin{cases}
  0 & (x<0)\\
  1 & (x>0)
 \end{cases}
\]
と定めれば,定義域のすべての点で $H'=0$ だが,\,$H$ は定数ではない.定義域が $x<0$ と $x>0$ の2つの区間に分かれているので,それぞれの区間で別の定数をとれるのである.
%
%
\section{増減と凹凸}\label{節:増減凹凸}
$f'(c)=0$ は極値の候補を探す条件であり,極値を保証する条件ではない.\,$f(x)=x^3$ は $f'(0)=0$ だが,\,$0$ の前後でも増加し続ける.極大・極小を判定するには,その点の前後で導関数の符号がどう変わるかを調べる.\,$+$ から $-$ に変われば極大,\,$-$ から $+$ に変われば極小である.

導関数をもう一度微分した $f''=(f')'$ を\textbf{第二次導関数}という.\,$f''>0$ なら接線の傾き $f'$ が増加し,グラフは\textbf{下に凸}になる.\,$f''<0$ なら傾きが減少し,グラフは\textbf{上に凸}になる.たとえば $y=x^2$ は下に凸で,\,$y=-x^2$ は上に凸である.

この呼び方を,接線との位置関係から確かめられる.\,$f''>0$ の区間で $a$ を固定し,\,$x>a$ とすると,平均値の定理により
\[
 f(x)-f(a)=f'(c)(x-a)>f'(a)(x-a)
 \qquad(a<c<x)
\]
となる.\,$x<a$ の場合も,\,$x<c<a$ で $f'(c)<f'(a)$ となることと $x-a<0$ を合わせれば,同じ結論を得る.したがって,曲線は接線の上側にある.\,$f''<0$ では不等号が逆になる.

グラフの凹凸が切り替わる点を\textbf{変曲点}という.\,$f''$ が連続な場合,凹凸が切り替わる場所は $f''=0$ となる点から探せるが,\,$f''=0$ だけでは十分でない.\,$y=x^4$ は $x=0$ で $f''=0$ になっても,その前後でともに下に凸である.

\begin{example}{極値と最大・最小}
$f(x)=x^3-3x$ の増減と凹凸を調べ,区間 $[-2,3]$ における最大値・最小値を求めよ.
\end{example}
導関数を因数分解すると,符号を調べやすい.
\[
 f'(x)=3x^2-3=3(x-1)(x+1),\qquad f''(x)=6x.
\]
したがって $x<-1$ で増加,\,$-1<x<1$ で減少,\,$x>1$ で増加する.\,$f(-1)=2$ が極大値,\,$f(1)=-2$ が極小値である.また,\,$x<0$ で上に凸,\,$x>0$ で下に凸だから,原点 $(0,0)$ は変曲点である.
\[
\begin{array}{c|ccccccc}
 x&-2&&-1&&1&&3\\ \hline
 f'(x)&&+&0&-&0&+&\\
 f(x)&-2&\nearrow&2&\searrow&-2&\nearrow&18
\end{array}
\]
閉区間での最大・最小は,極値だけを比較しては決まらない.両端も候補になるため,\,$f(-2)=-2$,\,$f(3)=18$ を加えて比較する.最大値は $18$ で $x=3$ のとき,最小値は $-2$ で $x=-2,1$ のときにとる.

\begin{figure}[htbp]
\centering
\begin{tikzpicture}[x=2.1cm,y=0.8cm,>=Stealth]
 \draw[->] (-2.3,0)--(2.35,0) node[right]{$x$};
 \draw[->] (0,-3.4)--(0,3.5) node[above]{$y$};
 \draw[very thick,domain=-2.1:2.1,samples=120] plot (\x,{\x^3-3*\x});
 \fill (-1,2) circle(1.5pt) node[above left]{$(-1,2)$};
 \fill (1,-2) circle(1.5pt) node[below right]{$(1,-2)$};
 \fill (0,0) circle(1.5pt) node[above right]{変曲点};
 \draw[dashed] (-1,0) node[below]{$-1$}--(-1,2)--(0,2);
 \draw[dashed] (1,0) node[above]{$1$}--(1,-2)--(0,-2);
 \node at (1.1,2.85) {$y=x^3-3x$};
\end{tikzpicture}
\caption{極値は増減の切り替わり,変曲点は凹凸の切り替わりを表す.}
\label{図:三次増減}
\end{figure}

図\ref{図:三次増減}は原点付近を描いたものである.見えている山の高さ $2$ が,指定区間全体の最大値とは限らないことに注意しよう.最大・最小を求めるときは,まず定義域と対象の区間を確かめ,内点の極値候補,微分できない点,含まれる端点を比較する.連続関数でも,区間が閉じていなければ最大・最小が存在しない場合がある.


%
\section{方程式と不等式}\label{節:微分方程式不等式}
方程式の根を式で表せなくても,根が存在することや,根の個数は分かることがある.ここでは,連続性が存在を保証し,増減が個数を制限する.不等式では,左右の差を一つの関数にして,その最小値を調べる.

\begin{example}{方程式と不等式}
\begin{enumerate}[label=(\arabic*),leftmargin=*]
 \item 方程式 $x=\cos x$ が区間 $[0,1]$ にただ一つの解をもつことを示せ.
 \item $x>0$ のとき $\log x\le x-1$ を示し,等号成立条件を求めよ.
\end{enumerate}
\end{example}
(1) では $h(x)=x-\cos x$ とおく.\,$h$ は $[0,1]$ で連続で,
\[
 h(0)=-1<0,\qquad h(1)=1-\cos1>0
\]
である.中間値の定理から,\,$h(x)=0$ となる点が $(0,1)$ に存在する.一方,この区間では $h'(x)=1+\sin x>0$ だから,\,$h$ は狭義に増加する.二つの異なる点でともに $0$ をとることはできず,解はただ一つである.グラフの交わりを予想しただけではなく,存在と一意性を別々に示した.

(2) は右辺から左辺を引き,\,$g(x)=x-1-\log x$ とおく.
\[
 g'(x)=1-\frac1x=\frac{x-1}{x}.
\]
定義域では $x>0$ なので,符号は分子 $x-1$ で決まる.\,$g$ は $0<x<1$ で減少し,\,$x>1$ で増加する.したがって,最小値は $g(1)=0$ であり,\,$g(x)\ge0$ を得る.両側で狭義に増減しているため,等号は $x=1$ のときに限る.


%
\section{一次近似}\label{節:一次近似}
\sectionref{節:微分定義}で得た式
\[
 f(a+h)=f(a)+f'(a)h+h\varepsilon(h),\qquad \varepsilon(h)\to0
\]
を使えば,\,$h$ が小さいとき
\[
 f(a+h)\approx f(a)+f'(a)h
\]
と近似できる.これを\textbf{一次近似}という.右辺は $a$ における接線の値である.計算しやすい点 $a$ を近くに選び,そこからの変化を傾きで見積もる.

\begin{example}{平方根の近似}
$\sqrt{4.04}$ を一次近似で求め,近似値が実際の値より大きいか小さいかを確かめよ.
\end{example}
$f(x)=\sqrt{x}$ とすると,\,$f(4)=2$,\,$f'(4)=\frac1{2\sqrt4}=\frac14$ である.\,$4.04=4+0.04$ なので,
\[
 \sqrt{4.04}\approx2+\frac14\cdot0.04=2.01.
\]
$f''(x)=-\cfrac1{4x^{\cfrac32}}<0$ だから,グラフは上に凸で,接線による値は実際の値より大きい.数値によっても,\,$2.01^2=4.0401>4.04$ と確かめられる.

誤差の大きさは,有理化によって直接評価できる.
\[
 2.01-\sqrt{4.04}
 =\frac{2.01^2-4.04}{2.01+\sqrt{4.04}}
 =\frac{0.0001}{2.01+\sqrt{4.04}}.
\]
分母は $4$ より大きいから,\,$0<2.01-\sqrt{4.04}<0.000025$ である.一次近似の式だけでは,指定した小数桁までの正しさが自動的に保証されるわけではない.この例では別の式で誤差を評価した.より一般的な誤差の扱いは\chapterref{章:近似と広義積分}で学ぶ.


%
\section{曲線の微分}\label{節:逆媒介陰微分}
\subsection{逆関数の微分}
一つの区間で連続かつ狭義単調な関数 $y=f(x)$ は,その値域で逆関数 $x=g(y)$ をもつ.このとき逆関数も連続であるという基本性質を用いる.その区間の内点 $a$ について $b=f(a)$,\,$f'(a)\ne0$ とし,\,$y=b+k$ に対応する $x$ を $a+h$ と書けば,\,$k\to0$ のとき $h\to0$ であり,\,$k=f(a+h)-f(a)$ である.したがって
\[
 \frac{g(b+k)-g(b)}k
 =\frac{h}{f(a+h)-f(a)}
 =\cfrac{1}{\cfrac{f(a+h)-f(a)}h}
 \longrightarrow\frac1{f'(a)}.
\]
よって
\[
 g'(b)=\frac1{f'(a)},\qquad
 g'(y)=\frac1{f'(g(y))}
\]
となる.分母が $0$ でないことが必要である.たとえば $f(x)=x^2$ を $x>0$ に限ると,逆関数は $g(y)=\sqrt y$ であり,\,$g'(y)=\frac1{2\sqrt y}$ となる.\,$x^2$ を実数全体で考えると,正負の二つの $x$ が同じ値に対応して逆関数が定まらないので,先にどちらの枝を使うかを決めなければならない.


%
\subsection{媒介変数表示}
曲線が $x=x(t)$,\,$y=y(t)$ と表されるとき,\,$t$ を\textbf{媒介変数}という.ここでは $x(t),y(t)$ が連続な導関数をもつとする.

考えている点で $x'(t)\ne0$ なら,\,$x'$ の連続性からその近くでも $x'(t)$ は同じ符号をもち,\,$x(t)$ は狭義単調になる.したがって,前の小節により $t$ を $x$ の関数 $t=t(x)$ として表せて,
\[
 \frac{\dd t}{\dd x}=\frac1{x'(t)}
\]
となる.\,$y=y(t(x))$ は $x$ の関数だから,合成関数の微分により
\begin{align*}
 \frac{\dd y}{\dd x}
 &=y'(t)\frac{\dd t}{\dd x}\\
 &=\frac{y'(t)}{x'(t)}
\end{align*}
を得る.\,$x'(t)=0$ の点ではこの商は使えず,垂直な接線などを別に調べる必要がある.

さらに2回微分できる場合は,同じ合成関数の微分から,\,$x$ で微分することが $\frac1{x'(t)}\frac{\dd}{\dd t}$ に対応する.よって $\frac{\dd y}{\dd x}$ をもう一度 $x$ で微分すると
\[
 \frac{\dd^2y}{\dd x^2}
 =\frac1{x'(t)}\frac{\dd}{\dd t}
   \left(\frac{y'(t)}{x'(t)}\right)
\]
となる.\,$\frac{y''(t)}{x''(t)}$ ではないことに注意しよう.
%
\subsection{陰関数の微分}
$x,y$ の関係が一つの式で与えられ,\,$y$ を $x$ の関数として明示していない場合がある.その式が実際に微分可能な関数 $y=y(x)$ を定めている範囲では,両辺を $x$ で微分できる.これを\textbf{陰関数の微分}という.\,$y$ も $x$ とともに変わるので,たとえば $(y^2)'=2yy'$,\,$(xy)'=y+xy'$ となる.

\begin{example}{接線の別解}
\begin{enumerate}[label=(\arabic*),leftmargin=*]
 \item $x=t^2+1$,\,$y=t^3-t$ で表される曲線の $t=1$ における接線を求めよ.
 \item 円 $x^2+y^2=25$ の点 $(3,4)$ における接線を求めよ.
\end{enumerate}
\end{example}
(1) では $t=1$ のとき $(x,y)=(2,0)$ である.また,
\[
 \frac{\dd x}{\dd t}=2t,\qquad
 \frac{\dd y}{\dd t}=3t^2-1,\qquad
 \left.\frac{\dd y}{\dd x}\right|_{t=1}
 =\frac{3-1}{2}=1.
\]
したがって接線は $y=x-2$ である.\,$t$ を消去しなくても,それぞれの変化率の比で傾きを求められた.

(2) では,点 $(3,4)$ の近くで上半円 $y=\sqrt{25-x^2}$ を考えれば,\,$y$ が微分可能な関数であることを確かめられる.その上で,元の関係式を微分すると
\[
 2x+2yy'=0,\qquad y'=-\frac xy\quad(y\ne0).
\]
点 $(3,4)$ での傾きは $-\frac34$ だから,接線は
\[
 y-4=-\frac34(x-3),\qquad 3x+4y=25
\]
である.円全体を一つの関数 $y=f(x)$ とみなしたわけではなく,指定点の近くの枝を使っている.\,$y=0$ の点では上の商は使えず,実際に接線は垂直になる.

接線を求める際には,曲線をどのように表すかによって計算が変わった.次は,最大値を求める場合にも,条件の表し方を変えて比べてみよう.
%
%
\section{幾何による別解}\label{節:幾何による別解}
微分を学ぶと,最大値や接線を同じ手順で求められるようになる.ただし,図形の性質から直接求まることもある.二つの方法を比べてみよう.

\begin{example}{円に内接する長方形}
半径$1$の円に内接する長方形の面積の最大値を求めよ.
\end{example}
円の中心を原点とし,長方形の辺が座標軸に平行になるように取る.第1象限の頂点を$(x,y)$とすれば,\,$x,y>0$,\,$x^2+y^2=1$であり,面積は$S=4xy$である.\,$y=\sqrt{1-x^2}$と表して微分すると,
\[
 S(x)=4x\sqrt{1-x^2},\qquad
 S'(x)=\frac{4(1-2x^2)}{\sqrt{1-x^2}}.
\]
$0<x<1$では分母は正なので,\,$x=\frac1{\sqrt2}$の前後で増加から減少へ変わる.このとき$y=\frac1{\sqrt2}$であり,最大値は$2$となる.

\noindent\textbf{【別解】}\quad
同じ式を$x,y$のまま見てみよう.平方が非負であることから,
\[
 (x-y)^2\ge0,\qquad 2xy\le x^2+y^2=1.
\]
したがって$S=4xy\le2$である.等号は$x=y=\frac1{\sqrt2}$のときに成り立つ.図形では,長方形が正方形になる場合である.

幾何的にも,長方形の辺を$a,b$,対角線を$2$とすれば,$a,b$を二辺とする直角三角形の斜辺は$2$である.この直角三角形を四つ並べ,外側を一辺$2$の正方形,中央を一辺$|a-b|$の正方形にすると,面積の和から
\[
 4=4\cdot\frac{ab}{2}+(a-b)^2
\]
となる.したがって$ab\le2$であり,等号は中央の正方形の面積が$0$,すなわち$a=b$のときに成り立つ.

微分による解法は,円を別の曲線へ変えても,一つの変数で面積を表せれば使える.一方,平方を使った解法では,円の条件$x^2+y^2=1$がそのまま上からの評価になる.答えの正方形も,等号条件$x=y$から分かる.このように,計算の長さだけでなく,条件のどの部分を使ったかまで比べておきたい.

%
\section{ニュートン法}\label{節:ニュートン法}
一次近似は関数の値を求めるだけでなく,方程式の解を探すためにも使える.\,$f(x)=0$の解に近そうな値を$x_n$とし,その点での接線を考えよう.接線は
\[
 y=f(x_n)+f'(x_n)(x-x_n)
\]
である.\,$f'(x_n)\ne0$なら,\,$x$軸との交点は
\[
 x=x_n-\frac{f(x_n)}{f'(x_n)}
\]
となる.この交点を次の値$x_{n+1}$として繰り返す方法を\textbf{ニュートン法}という.曲線上の点から接線を引き,接線の零点から再び曲線へ戻る計算である.

\begin{example}{平方根に近づく数列}
$x_1=2$とし,\,$f(x)=x^2-2$にニュートン法を使って得られる数列を求めよ.また,その数列が$\sqrt2$に収束することを示せ.
\end{example}
$f'(x)=2x$だから,
\[
 x_{n+1}=x_n-\frac{x_n^2-2}{2x_n}
 =\frac12\left(x_n+\frac2{x_n}\right).
\]
順に計算すると,\,$x_2=\frac32$,\,$x_3=\frac{17}{12}$となる.ただし,値が近づいて見えるだけでは収束の証明にはならない.

まず,\,$x_n>\sqrt2$なら,
\[
 x_{n+1}-\sqrt2
 =\frac{x_n^2-2\sqrt2x_n+2}{2x_n}
 =\frac{(x_n-\sqrt2)^2}{2x_n}>0
\]
となる.\,$x_1>\sqrt2$なので,すべての項は$\sqrt2$より大きい.さらに,
\[
 x_{n+1}-x_n=\frac{2-x_n^2}{2x_n}<0
\]
だから単調減少する.下に有界なので極限$L\ge\sqrt2$をもち,漸化式の両辺の極限を取れば
\[
 L=\frac12\left(L+\frac2L\right),\qquad L^2=2
\]
である.\,$L>0$より$L=\sqrt2$と分かる.

途中の式は,誤差の小さくなる様子も表している.\,$e_n=x_n-\sqrt2$とおくと,
\[
 e_{n+1}=\frac{e_n^2}{2x_n}.
\]
例えば$0<e_n<\frac1{10}$なら,\,$x_n>1$なので$e_{n+1}<\frac1{200}$である.接線から作った数列を,差や大小関係を調べる漸化式の問題として考え直せた.

ただし,どの関数でもこのように収束するわけではない.例えば$f(x)=x^3-2x+2$で$x_1=0$とすると,\,$x_2=1$,\,$x_3=0$となり,二つの値を繰り返す.接線の零点を使う理由と,実際に収束することは,それぞれ確かめる必要がある.

%
\section{微分方程式}\label{節:微分方程式}
これまでは,関数が与えられたときに導関数を求めた.逆に,導関数についての条件から関数を決めることもできる.例えば
\[
 y'=ky
\]
を満たす関数$y$を考えよう.\,$k$は定数である.未知数が数ではなく関数であり,その導関数も含む方程式を\textbf{微分方程式}という.

$y=\e^{kx}$ならこの式を満たす.定数倍した$C\e^{kx}$も満たすが,ほかの形の解はないだろうか.そのことを調べるため,$y$を既に見つけた解$\e^{kx}$で割ってみよう.\,$y=C\e^{kx}$なら,この商は定数$C$になるはずである.解$y$に対して実際に微分すると,
\[
 \frac{\dd}{\dd x}(\e^{-kx}y)
 =\e^{-kx}(y'-ky)=0
\]
となる.平均値の定理から,導関数が$0$の関数は一つの区間で定数である.したがって$\e^{-kx}y=C$より$y=C\e^{kx}$となる.\,$y(0)=A$も与えられれば$C=A$であり,解は一つに定まる.

等しい間隔$h$での値だけを取り出して$a_n=y(nh)$とおけば,
\[
 a_{n+1}=\e^{kh}a_n
\]
である.この条件から見れば等比数列となっている.微分方程式を学ぶ前から知っていた等比的な増加も,間の値を指数関数で表すことができる.

もう一つ,\,$y''=-y$を考えてみよう.\,$\sin x,\cos x$はいずれも二回微分すると符号が変わるので,
\[
 y=A\cos x+B\sin x
\]
は解である.さらに$y(0)=A$,\,$y'(0)=B$を満たす.これ以外の解がないことも確かめられる.同じ初期条件をもつ別の解との差を$z$とすると,
\[
 z''=-z,\qquad z(0)=z'(0)=0
\]
である.\,$z''=-z$の両辺に$z'$を掛けると,$z'z''=-zz'$となる.左辺は$\frac12(z')^2$の微分,右辺は$-\frac12z^2$の微分だから,二つを同じ側に集めた$z^2+(z')^2$を調べればよい.実際,
\[
 \{z^2+(z')^2\}'=2zz'+2z'z''=2z'(z+z'')=0.
\]
したがってこの式は定数で,\,$x=0$で$0$だから常に$0$となる.二つの平方の和が$0$なので$z=0$である.微分方程式の一般的な理論を使わずに,解の形と一意性を確かめられた.

時刻$t$での位置を$s(t)$とすると,速度は$s'(t)$であり,速度の変化率$s''(t)$を\textbf{加速度}という.例えば$s''(t)=-s(t)$は,原点から離れるほど大きく,原点へ向かう加速度が生じる運動を表す.先ほどの結果から,位置は$s(t)=A\cos t+B\sin t$となり,周期的に往復する.このような運動を\textbf{単振動}という.同じ微分方程式は,変数を時刻と考えない場合にも,二回微分したときの関係によって三角関数を特徴付けている.

ここまでは,微分から分かることを調べた.次章では,導入で扱った面積を一般の関数について定め,その変化率が元の関数になることを示す.そこから,微分の逆によって積分を計算する方法を考えよう.
%
\section*{章末問題}
\begin{enumerate}[label=\textbf{\thechapter.\arabic*},leftmargin=*]

 \item\label{問:平均値評価} 平均値の定理を使い,任意の実数 $a,b$ に対して $|\sin b-\sin a|\le|b-a|$ が成り立つことを示せ.適用する関数と区間を明記せよ.
 \item\label{問:三次関数グラフ} $f(x)=x^3-3x^2$ の増減,極値,凹凸,変曲点を求めよ.さらに,区間 $[-1,3]$ における最大値・最小値と,それらをとる $x$ を求めよ.
 \item\label{問:指数不等式方程式} 微分を用いて,すべての実数 $x$ について $\e^x\ge1+x$ を示し,等号成立条件を求めよ.また,方程式 $\log x+x=2$ は $x>0$ でただ1つの解をもち,その解が $(1,2)$ にあることを示せ.
 \item\label{問:逆媒介陰} 次を求めよ.
 \begin{enumerate}[label=(\arabic*),leftmargin=*]
  \item $f(x)=x^3+x$ の逆関数を $g$ とするとき,\,$g'(0)$.
  \item $x=t^2$,\,$y=t^3$ で表される曲線の,\,$t=2$ における接線.
  \item 曲線 $x^2+xy+y^2=7$ の点 $(1,2)$ における接線.ただし,その点の近くで $y$ は $x$ の微分可能な関数として表されるものとする.
 \end{enumerate}

 \item\label{問:微分方程式初期値} 1つの区間で微分可能な関数$y$が$y'=3y$,\,$y(0)=2$を満たすとする.\,$\e^{-3x}y$を微分することにより$y$を求めよ.

 \item\label{問:三平方根ニュートン} $x^2-3=0$にニュートン法を用いた漸化式を求めよ.\,$x_1=2$とするとき,この数列が単調減少し,\,$\sqrt3$に収束することを示せ.
\end{enumerate}
\smallskip\noindent 解答は\pageref{解答:第3章}ページ.
%
